% State-specific — Compound Interest Introduction
\section{Compound Interest Introduction}

\begin{learningGoals}
\begin{itemize}
    \item Understand the difference between simple and compound interest
    \item Calculate compound interest for short periods
\end{itemize}
\end{learningGoals}

\begin{conceptBox}{Simple vs.\ Compound Interest}
\textbf{Simple interest} is calculated on the \textbf{original principal only}.

\textbf{Compound interest} is calculated on the principal \textbf{plus} any interest already earned. Your interest earns interest!

$$A = P(1 + r)^t$$
\begin{itemize}[leftmargin=*]
    \item $A$ = total amount after $t$ years
    \item $P$ = principal (starting amount)
    \item $r$ = annual interest rate (as a decimal)
    \item $t$ = number of years
\end{itemize}
\end{conceptBox}

\begin{workedExample}{Comparing Simple and Compound Interest}
Deposit $\$500$ at $10\%$ for $3$ years.

\textbf{Simple:} $I = 500 \times 0.10 \times 3 = \$150$. Total $= \$650$.

\textbf{Compound (yearly):}
Year 1: $500 \times 1.10 = \$550$
Year 2: $550 \times 1.10 = \$605$
Year 3: $605 \times 1.10 = \$665.50$
\answerHighlight{Compound earns $\$15.50$ more than simple.}
\end{workedExample}

\mascotSays{Compound interest is like a snowball --- it starts small but grows faster and faster!}

\begin{practiceBox}{Compound Interest}
\resetProblems

\prob You deposit $\$1{,}000$ at $5\%$ compounded yearly. Find the total after $2$ years.
\answerExplain{$\$1{,}102.50$}{Year 1: $1{,}000 \times 1.05 = \$1{,}050$. Year 2: $1{,}050 \times 1.05 = \$1{,}102.50$.}

\prob $\$200$ at $10\%$ simple interest for $2$ years vs.\ compound. Find both totals.
\answerExplain{Simple: $\$240$; Compound: $\$242$}{Simple: $I = 200 \times 0.10 \times 2 = \$40$; total $= \$240$. Compound: $200 \times 1.10 = 220$; $220 \times 1.10 = \$242$.}

\prob Use $A = P(1+r)^t$ to find the value of $\$800$ at $4\%$ for $3$ years.
\answerExplain{$\$899.89$}{$A = 800(1.04)^3 = 800 \times 1.124864 = \$899.89$ (rounded to the nearest cent).}

\prob Why does compound interest grow faster than simple interest?
\answerExplain{You earn interest on your interest}{With compound interest, each year the base grows because previous interest is added. Simple interest always uses the original principal.}

\prob A $\$600$ loan at $8\%$ compounds yearly for $2$ years. How much is owed?
\answerExplain{$\$699.84$}{Year 1: $600 \times 1.08 = \$648$. Year 2: $648 \times 1.08 = \$699.84$.}

\end{practiceBox}
